\documentclass[11pt,a4paper]{article}
\usepackage[margin=2.1cm]{geometry}\usepackage{fontspec}\usepackage[italian]{babel}\usepackage{xcolor,longtable,booktabs,hyperref,fancyhdr}
\setmainfont{NimbusSans-Regular.otf}[Path=sorgenti/font/,BoldFont=NimbusSans-Bold.otf,ItalicFont=NimbusSans-Italic.otf]
\setmonofont{lmmono10-regular.otf}[Path=sorgenti/font/,Scale=.88]
\definecolor{teal}{HTML}{087C74}\hypersetup{colorlinks=true,linkcolor=teal,urlcolor=teal}\setlength{\headheight}{14pt}\setlength{\parindent}{0pt}\setlength{\parskip}{7pt}\pagestyle{fancy}\fancyhf{}\fancyfoot[R]{\thepage}\fancyhead[L]{Geologia · Statistica e coding R · R6}\begin{document}


\section*{Guida alla revisione integrale R6}

Il corso comprende 17 moduli, da T00 a M16, e 11 macroprogetti, da P01 a P11, per 861 slide. La revisione riunisce le sequenze ripetitive del materiale iniziale in pagine argomentate, conservando il raccordo con tutte le 2150 slide originali. Il registro è accademico e la prosa esplicita relazioni logiche, ipotesi e conseguenze dei calcoli.

\subsection*{Materiali della consegna}

Sono inclusi 261 grafici PDF prodotti in R, 232 diagrammi TikZ utilizzati nelle slide e 441 slide con formalismo matematico esplicito nei riquadri. I 340 laboratori R e i 108 estratti delle slide sono stati eseguiti con esito positivo; il percorso introduttivo conserva inoltre 20 laboratori Python verificati.

Il template delle slide mantiene formato 20 × 12,5 cm, caratteri Nimbus Sans, codice monospaziato e palette del materiale approvato. Le figure utilizzano assi leggibili e file vettoriali. I diagrammi sono sorgenti TikZ modificabili e compilati nelle slide.

\subsection*{Come leggere i materiali}

Aprire Indice\_corso.html dopo aver estratto lo ZIP. Ogni titolo porta al PDF corrispondente; gli indici interni dei fascicoli consentono la navigazione fra sezioni. Le cartelle R, tikz ed estratti contengono i sorgenti delle figure, dei diagrammi e dei frammenti mostrati. Il catalogo dei laboratori associa identificativo, titolo e file eseguibile.

\subsection*{Rilettura slide per slide}

La revisione R6 ha riesaminato tutte le 861 pagine, registrando un’osservazione individuale e intervenendo su 357 slide. Il controllo comprende autonomia delle spiegazioni, definizioni dei simboli, rapporto fra testo e rappresentazione, leggibilità di assi ed etichette e composizione nel template originale. Le pagine corrette sono state nuovamente compilate e controllate visivamente.

Il rapporto e il registro completo si trovano nella cartella audit. Il registro HTML consente di cercare modulo, pagina e osservazioni; i dati sono disponibili anche in CSV e JSON. L’autonomia è valutata rispetto ai prerequisiti del modulo: la singola pagina esplicita il proprio ragionamento senza ricostruire ogni nozione dell’intero corso.

Fra gli interventi sostanziali figurano la correzione dell’ombreggiatura del terreno in M12, il controllo della geometria dei diagrammi, le definizioni locali nelle formule e il contrasto dei numeri dei blocchi in P07. Il codice R modificato è stato rieseguito; il laboratorio M12 N12 è stato ricollaudato e il suo hash aggiornato dopo l’esecuzione positiva.

\subsection*{Interventi editoriali}

Le aperture ridondanti, le successioni ripetute di fase–codice–risultato e le sintesi sovrapposte sono state integrate nel percorso. Il controllo lessicale ha esaminato formule retoriche, contrapposizioni stereotipate e frasi introduttive generiche, conservando le negazioni che esprimono limiti scientifici necessari. Il raccordo slide per slide documenta gli accorpamenti: non vi è corrispondenza uno a uno fra numerazione originale e revisione.

\newpage\section*{Indice dei moduli}

\begin{longtable}{p{1.2cm}p{12cm}r}\toprule Codice & Titolo & Slide\\\midrule\endhead

T00 & Dal codice al calcolo scientifico & 33\\

M01 & Pensare con i dati & 58\\

M02 & R e gestione dei dati & 43\\

M03 & Descrivere e visualizzare i dati & 33\\

M04 & Probabilità, variabilità e simulazione & 35\\

M05 & Campionamento, stima e incertezza & 105\\

M06 & Confrontare gruppi e valutare ipotesi & 35\\

M07 & Relazioni tra variabili e regressione & 33\\

M08 & Analisi multivariata dei campioni & 33\\

M09 & Serie temporali ambientali & 33\\

M10 & Previsione e progetto statistico & 33\\

M11 & Statistica spaziale & 32\\

M12 & Raster, telerilevamento e morfometria & 30\\

M13 & Cronologie e archivi sedimentari & 30\\

M14 & Estremi, pericolosità e rischio & 31\\

M15 & Riproducibilità e progetto di analisi & 27\\

M16 & Classificazione e validazione & 29\\

\bottomrule\end{longtable}\newpage\section*{Indice dei macroprogetti}

\begin{longtable}{p{1.2cm}p{12cm}r}\toprule Codice & Titolo & Slide\\\midrule\endhead

P01 & Geometria dei pori e permeabilità & 21\\

P02 & Foraminiferi, coppe e campagne & 22\\

P03 & Il Nilo: livelli e dipendenza temporale & 21\\

P04 & Un geyser: associazione e previsione & 21\\

P05 & CO\textsubscript{2}: tendenza e calendario stagionale & 22\\

P06 & Dai campioni alla validazione spaziale & 17\\

P07 & Terreno, pendenze e sintesi territoriali & 17\\

P08 & Collocare nel tempo il massimo di un proxy & 18\\

P09 & Un livello raro e la sua incertezza & 16\\

P10 & Un’analisi ricostruibile e verificabile & 16\\

P11 & Classificare nuove stazioni sintetiche & 17\\

\bottomrule\end{longtable}

I macroprogetti integrano le attività originarie in analisi continue. Le nuove numerosità Monte Carlo sono dichiarate negli script e nelle slide; i laboratori conservati mantengono le proprie impostazioni. Il riuso didattico degli stessi dati viene esplicitato, soprattutto nella validazione predittiva.

\newpage\section*{Ricostruire analisi e presentazioni}

Aprire Corso\_Geologia.Rproj oppure impostare come cartella corrente la radice del progetto. Per la sola lettura dei PDF non occorre installare R, Python o LaTeX.

\subsection*{Ambiente R}

Il collaudo utilizza R 4.5 e ggplot2 3.5.1; le versioni effettive dei pacchetti sono registrate in verifiche/versioni\_R.csv e la sessione in verifiche/ambiente\_R.txt. Installare esplicitamente le dipendenze prima dell’analisi tramite R/00\_installazione.R, quindi eseguire R/controlla\_ambiente.R. sf e terra possono richiedere GDAL, GEOS, PROJ e UDUNITS. Il macroprogetto P10 richiede Git.

Per riprodurre il carattere delle figure Cairo, installare Nimbus Sans sul sistema operativo usando i font inclusi in sorgenti/font. XeLaTeX legge direttamente i font inclusi; la grafica Cairo di R utilizza invece il catalogo dei font del sistema.

\begin{verbatim}source("R/M07_figure.R")
source("laboratori/M07/M07_N01.R")
source("estratti/P04/separazione.R")
\end{verbatim}

I nomi degli estratti disponibili si leggono nella cartella estratti; il comando precedente illustra la struttura del percorso.

\subsection*{Ricostruzione automatizzata}

Python richiede i pacchetti elencati in strumenti/requirements.txt. La compilazione richiede XeLaTeX con Beamer, TikZ, fontspec, listings, bookmark e i pacchetti del template. Il programma seguente supporta un percorso R personalizzato con --r-bin.

\begin{verbatim}python strumenti/ricostruisci.py figure M07
python strumenti/ricostruisci.py slide M07
python strumenti/ricostruisci.py laboratori
python strumenti/ricostruisci.py estratti
python strumenti/ricostruisci.py python
python strumenti/ricostruisci.py tutto
\end{verbatim}

\newpage\section*{Verifiche e tracciabilità}

La compilazione controlla altezza di titoli, testo, figure e richiami, caratteri mancanti e righe eccedenti. Tutte le pagine sono state esaminate in tavole di contatto, con ingrandimenti mirati. Un controllo aggiuntivo verifica che il testo delle figure PDF rimanga entro la pagina. I risultati sono conservati nella cartella verifiche.

Il confronto indipendente delle covarianze di P03, P04 e P05 utilizza statsmodels; gli scarti assoluti massimi sono dell’ordine di 10 elevato a −12 o inferiori. Le dieci medie ricampionate condivise fra R e Python in M15 e P10 differiscono al massimo di circa 3,6 × 10 elevato a −15. Le tolleranze e le procedure sono esplicite nei relativi script.

Il catalogo dei laboratori conserva gli hash delle versioni eseguite. MANIFEST\_SHA256.txt permette di verificare l’integrità dei file consegnati. La rigenerazione di risultati, log o PDF può modificare gli hash: il manifesto descrive la consegna, mentre i controlli numerici verificano la ricostruzione del calcolo.

\subsection*{Provenienza e limiti}

I dataset sintetici sono dichiarati come tali. I casi rock, faithful, Nile, CO\textsubscript{2} e foraminiferi conservano le rispettive fonti e condizioni di provenienza; il dizionario e le note dei dati accompagnano il progetto. L’inferenza è delimitata dalle unità osservate, dalla dipendenza e dalle ipotesi del modello, evitando di attribuire ai risultati didattici una validità territoriale o operativa non verificata.

Nella cartella documentazione, fonti\_online\_R6.json raccoglie i riferimenti consultati e fonti\_M05.txt conserva la bibliografia specifica del modulo di stima. Le figure sono generate dai codici del corso e non riprodotte dai manuali. Ulteriori riferimenti: \path{verifiche/fonti_aggiuntive_R6.md}.

\end{document}